\section{Discussion}
\label{sec:discussion}

\paragraph{Multi-signal supervision versus simple heuristics.} The controlled results
show the mechanisms working as specified, while leaving open whether the multi-worker
design adds value over transparent rules: fixed-threshold and MAD rules tie the
fast-closure detector, a fastest-median-closure ranking equals or nearly equals SAT-SA
in most generated populations---populations that favour SAT-SA by construction---and
four workers \emph{lowered} population recall in the ablation. This is consistent with
a generator whose pathological profile includes fast closures; more elaborate workers
can only earn their place on data where pathologies are not separable by one rate,
which the present evaluation does not contain. For a supervisor, a small set of transparent
rates may be a reasonable first instrument, with record-cited findings serving as the
explanation layer rather than the ranking signal.

\paragraph{Mechanism validation versus external validity.} Controlled experiments with
construction labels answer whether each mechanism does what it declares; they cannot
answer whether its output reflects supervisory quality. The external evaluation shows
why the distinction matters: the same fast-closure signal that is a warning in a SOC is
a success indicator in a service desk, so a construct defined for one setting inverts
in the other. The methodological lesson generalizes beyond SAT-SA: an external outcome
used to ``validate'' a supervisory score must measure the same construct, and a
negative association with a different construct is evidence about constructs, not a
defect to tune away. The external study also exposed a design limitation within SAT-SA's
own domain---existence rules fused by summed confidence saturate on large
submissions---which calls for prevalence-aware scoring developed on held-out SOC data.

\paragraph{Orchestration versus overhead.} On this workload, checkpointed orchestration
cost a median \V{O01-diff-proc}\,s per workflow with an interval that includes zero, at
the price of \V{O01-diff-db} additional database statements; in exchange, all evaluated
interruptions recovered without duplicated decisions. Whether this trade holds under
load, on PostgreSQL or across machines was not measured.

\paragraph{Integrity versus cost.} Finalization and verification are sub-second on this
machine, and the tested alterations of the decided record were detected. The guarantee
is narrow by construction: it binds what was analysed and decided, and it is only as
strong as key custody, which was not evaluated.

\paragraph{Peer analysis versus cohort dependence.} Peer comparison is the most
supervisor-specific signal, but P01 shows that its output is a function of cohort size
and spread as much as of the subject. Cohort definition is therefore a supervisory
decision that should be recorded and reviewed with the finding.

\paragraph{The human in the loop.} SAT-SA is decision support. Automation bias is a
documented risk for such tools~\cite{parasuraman2010automation,tilbury2024automationbias};
without a study with supervisors we cannot say whether record-cited findings help or
anchor their judgement.
